\section{Commutative Algebra}

\subsection{Rings and Ideals}

We recall some basic notions regarding rings.

\bdefn

    A {\emphcolor ring} is a set $R$ equipped with two operations $+$ and $\cdot$, making it into a commutative group and a monoid respectively.
    Furthermore, they satisfy the distributivity properties:
    $$ a\cdot(b+c) = ab + ac,\qquad (a+b)\cdot c = ac + bc $$
    for all $a,b,c\in R$.
    The additive and multiplicative identities are denoted $0$ and $1$ respectively.

\edefn

\bnote

    In this course all rings are assumed to be {\emphcolor commutative}: $\cdot$ forms a commutative monoid.

\enote

\bdefn

    A {\emphcolor ring morphism} is a function between two rings, $f\colon R\to S$, which preserves structure: $f(x+y)=f(x)+f(y)$, $f(xy)=f(x)f(y)$,
    and $f(1)=1$.
    A {\emphcolor ring isomorphism} is a ring morphism with an inverse (equivalently a bijective ring morphism).

    A {\emphcolor subring} of a ring $R$ is a subset $S\subseteq R$ equipped with a ring structure making the inclusion $S\to R$ a ring morphism.

    An {\emphcolor ideal} of a ring $R$ is a nonempty subset $I\subseteq R$ closed under addition and multiplication by $R$: for $r\in R$ and $x\in I$,
    $rx\in I$.
    For an ideal $I\subseteq R$, we define its {\emphcolor quotient ring} $R/I$ to be the set of all cosets $[r]=r+I$ with ring structure given by
    $$ 0 = [0] = I,\qquad 1 = [1],\qquad [x] + [y] = [x+y],\qquad [x]\cdot[y] = [xy] , $$
    equivalently the unique structure making the projection map $\pi\colon R\to R/I,r\mapsto[r]$ a ring morphism.

\edefn

\bprop

    Let $f\colon R\to S$ be a ring morphism, then
    \benum
        \item For every ideal $J\leq S$, $f^{-1}(J)\leq R$ is an ideal.
        \item In particular $f$'s {\emphcolor kernel}, $\ker f=f^{-1}0$, is an ideal of $R$.
        And $f$'s {\emphcolor image}, $\im f=f(R)$ is a subring of $S$.
        $f$ then quotients to give a ring isomorphism $R/\ker f\cong\im f$.
        \item For an ideal $I\leq R$, consider the projection $\pi\colon R\to R/I$.
        Then $J\mapsto\bar J=\pi(J)$ and $\bar J\mapsto\pi^{-1}\bar J$ gives a one-to-one correspondence between ideals of $R$ containing $I$
        and ideals of $R/I$.
    \eenum

\eprop

\bdefn

    For a ring $R$ and a subset $S\subseteq R$, we denote the smallest ideal containing $S$ by $(S)$.
    This exists as the arbitrary intersection of ideals is an ideal.
    When $S=\set x$ is a singleton, we denote $(\set x)=(x)$; such an ideal is called a {\emphcolor principal ideal}.

\edefn

\bprop

    The ideal generated by $S$, $(S)$, is the set of all finite $R$-combinations of elements from $S$:
    $$ (S) = \set{\sum_{s\in S}r_ss}[r_s\in R,\hbox{ all but finitely many being zero}] . $$
    In particular, $(x)=Rx$ is the set of all multiples of $x$.

\eprop

\bproof

    It is not hard to see that this is an ideal, and clearly every ideal containing $S$ must contain it.
    \qed

\eproof

\bdefn

    Let $R$ be a ring, an element $x\in R$ is
    \benum
        \item a {\emphcolor zero divisor} if there is a $y\in R$ such that $xy=0$;
        \item a {\emphcolor unit} if there is a $y\in R$ such that $xy=1$;
        \item {\emphcolor nilpotent} if $x^n=1$ for some $n>0$.
    \eenum
    The set of units of $R$ is denoted $R^\times$ and forms a group.

\edefn

\bdefn

    A ring $R\neq0$ is
    \benum
        \item {\emphcolor reduced} if every $x\neq0$ is not nilpotent;
        \item an {\emphcolor integral domain} if every $x\neq0$ is not a zero divisor;
        \item a {\emphcolor field} if every $x\neq0$ is a unit.
    \eenum

\edefn

It is easy to see that a field is an integral domain is reduced.

\bprop

    For a ring $R$, the following are equivalent:
    \benum
        \item $R$ is a field,
        \item $R$ has no non-trivial ideals,
        \item Every ring morphism out of $R$ to a non-zero ring is injective.
    \eenum

\eprop

\bproof

    If $R$ is a field, then every nonzero element is a unit and thus the ideal generating it must be all of $R$.
    Conversely if $R$ has no non-trivial ideals, then for $x\neq0$ we must have that $1\in(x)$ and thus there is a $y\in R$ such that $xy=1$ so it is a unit.

    If $R$ has no non-trivial ideals then for every $f\colon R\to S\neq0$ we must have $\ker f=0$, as it is an ideal not containing $1$.
    Thus $f$ is injective.
    Conversely if $I\neq R$ is an ideal of $R$, then $\pi\colon R\to R/I$ is injective, but its kernel is $I$, so $I=0$.
    \qed

\eproof

\bprop

    If $\set{I_\alpha}_\alpha$ is a collection of ideals from $R$, then
    \benum
        \item $\bigcap_\alpha I_\alpha$ is an ideal, the largest ideal contained in each $I_\alpha$.
        \item $\sum_\alpha I_\alpha=\set{\sum_\alpha x_\alpha}[x_\alpha\in I_\alpha\hbox{ all but finitely many are zero}]$ is an ideal, the smallest ideal containing each $I_\alpha$.
    \eenum

\eprop

\bproof

    The first point is easily verified, and the second one follows because
    $$ \sum_\alpha I_\alpha = \parens{\bigcup_\alpha I_\alpha} . \qed $$

\eproof

\bdefn

    If $I,J$ are ideals then we define their {\emphcolor product} to be the ideal generated from all products of elements from $I$ and $J$:
    $$ IJ = (\set{xy}[x\in I,y\in J]) . $$

\edefn

Since $I,J$ are ideals and are thus closed under multiplication from $R$, for $x\in I$ and $y\in J$ we have that $xy\in I\cap J$.
Thus $IJ\subseteq I\cap J$.

\bdefn

    For an ideal $I\subseteq R$, we define its {\emphcolor radical} to be
    $$ \sqrt I = \set{r\in R}[r^n\in I\hbox{ for some $n\geq0$}] . $$

\edefn

\blemm

    The radical of an ideal is an ideal.

\elemm

\bproof

    If $a^n\in I$ and $r\in R$ then $(ra)^n=r^na^n\in I$ so that $ra\in\sqrt I$.
    And if $a^n,b^m\in I$ then
    $$ (a+b)^{n+m} = \sum_{i=0}^{n+m}\binom{n+m}i a^ib^{n+m-i} \in I , $$
    since for each $i$, if $i<n$ then $n+m-i>m$, so $a+b\in\sqrt I$.
    \qed

\eproof

\subsection{Prime and maximal ideals}

\bdefn

    A proper ideal $I\subset R$ is
    \benum
        \item {\emphcolor prime} if $ab\in I$ implies $a\in I$ or $b\in I$.
        Equivalently if $a,b\notin I$ then $ab\notin I$.
        \item {\emphcolor maximal} if it is not contained in any other proper ideal.
    \eenum

\edefn

\bprop

    An ideal $I$ is
    \benum
        \item prime iff $R/I$ is an integral domain;
        \item maximal iff $R/I$ is a field.
    \eenum
    Consequently, a maximal ideal is prime.

\eprop

\bproof

    $R/I$ is an integral domain iff $[a][b]\neq[0]$ when $[a],[b]\neq[0]$.
    This means that $ab\notin I$ when $a,b\notin I$, as desired.

    $R/I$ is a field iff it has no nontrivial ideals.
    Since ideals of $R/I$ are in bijection with ideals containing $I$, this is iff $I$ is maximal.
    \qed

\eproof

In particular the zero ideal $0=(0)=\set0$ is prime iff $R\cong R/(0)$ is an integral domain.

\bcoro

    Let $I$ be an ideal, then the correspondence $J\oto\bar J$ between ideals containing $I$ and ideals of $R/I$
    preserves primality and maximality: $J$ is prime/maximal iff $\bar J$ is.

\ecoro

\bproof

    We have that
    $$ (R/I)/\bar J = (R/I)/(J/I) \cong R/J , $$
    so that $(R/I)/\bar J$ is an integral domain/field iff $R/J$ is.
    The result follows.
    \qed

\eproof

\bdefn

    \benum
        \item An element $0\neq x\in R$ is {\emphcolor irreducible} if it is not a unit, and if $x=yz$ then either $y$ or $z$ is a unit.
        \item An integral domain $R$ is a {\emphcolor principal ideal domain} (or a {\emphcolor PID}) if all of its ideals are principal.
    \eenum

\edefn

\bprop

    If $R$ is a PID and $0\neq x$, then $(x)$ is prime iff $x$ is irreducible iff $(x)$ is maximal.

\eprop

\bproof

    If $(x)$ is prime, $x$ is not a unit since it is a proper ideal.
    And if $x=yz$ we have that $y$ or $z$ is in $(x)$, say $y=xt$.
    Then $x=yz=xtz$, and so $x(1-tz)=0$, and since $R$ is an integral domain this means $tz=1$ so $z$ is a unit.

    Now if $x$ is irreducible, then suppose $(x)\subseteq I$ is a proper ideal containing $(x)$.
    Since $R$ is a PID, $I=(y)$ so that $x\in(y)$ meaning $x=yz$.
    $(y)$ is a proper ideal so $y$ is not a unit, thus $z$ is, and then $y=xz^{-1}$ so that $y\in(x)$ and thus $(x)=(y)$.
    So $(x)$ is maximal, as desired.
    \qed

\eproof

\bexam

    The ring of integers $\bZ$ is a PID.
    Its irreducible elements are precisely its prime elements, so that its prime ideals are $(0)$ and $(p)$ for $p\in\bZ$ prime.

\eexam

\bexam

    If $k$ is a field, then its ring of polynomials $k[x]$ is a PID.
    Indeed, if $0\neq I\leq k[x]$ is an ideal, then let $f\in I$ have minimal degree.
    By polynomial division we see that $I=(f)$: for $g\in I$ we have that $g=fq+r$, but since $r=g-fq\in I$ has smaller degree than $f$,
    it must be zero.

    Its prime ideals are $(0)$ and $(f)$ for an irreducible monic polynomial $f$.

\eexam

\blemm

    If $f\colon R\to S$ is a ring morphism, and $\q\leq S$ is prime, then $f^{-1}\q\leq R$ is prime.

\elemm

\bproof

    A direct proof is easy, instead consider the composition:
    $$ \bar f\colon R\xtos fS\xtos\pi S/\q . $$
    By definition $f^{-1}\q=\ker\bar f$, and so $\bar f$ induces an embedding $R/f^{-1}\q\inj S/\q$.
    Subrings of integral domains are integral domains, so $R/f^{-1}\q$ is an integral domain, thus $f^{-1}\q$ is prime.
    \qed

\eproof

Subrings of fields are not fields, so the above proof does not show that the preimage of a maximal ideal is maximal.
Indeed, this is not true: consider the unique map $f\colon\bZ\to\bQ$.
Then $0\leq\bQ$ is maximal, but $f^{-1}0=\ker f=0$ is not maximal in $\bZ$.

\bthrm

    For every ring $R$ and every proper ideal $I<R$, there is a maximal ideal $\m$ containing $I$.

\ethrm

\bproof

    Consider the set of all proper ideals containing $I$.
    This is nonm-empty, as it contains $I$.
    Furthermore, the union of every chain of proper ideals is a proper ideal; it being an ideal is easy to see, and it is proper
    since $1$ is not in any element of the chain.
    Thus by Zorn, this set has a maximal element, which must be a maximal ideal.
    \qed

\eproof

\bcoro

    Let $x\in R$, then $x$ is a unit iff $x$ is not in any maximal ideal.
    That is,
    $$ R^\times = \bigcap_\m\m^c , $$
    where the intersection is over all maximal ideals of $R$.

\ecoro

\bproof

    An element is a unit iff it is not contained in any proper ideal.
    This is clear enough: if $x$ is not a unit then $(x)$ is proper, and proper ideals don't contain units.
    Every proper ideal can be extended to a maximal ideal, so the result follows.
    \qed

\eproof

\bdefn

    Let $R$ be a ring.
    Then define its {\emphcolor nilradical} to be $\nil(R)=\sqrt0$, i.e.\ the ideal of all nilpotent elements.

\edefn

Clearly $R$ is reduced iff $\nil(R)=0$.

\bprop

    Let $I\leq R$ be an ideal, and let $\pi\colon R\to R/I$ be the projection.
    Then $\sqrt I=\pi^{-1}\nil(R/I)$.

\eprop

\bproof

    We note that if $x^n\in I$ then $\pi(x)^n\in I$ so that $\pi(x)\in\nil(R/I)$ as desired.
    Conversely, if $\pi(x)\in\nil(R/I)$ then $\pi(x^n)=I$ so that $x^n\in I$ as desired.
    \qed

\eproof

\bprop

    For a ring $R$,
    \benum
        \item $R/\nil(R)$ is reduced;
        \item The nilradical is the intersection of all prime ideals:
        $$ \nil(R) = \bigcap_\p\p . $$
    \eenum

\eprop

\bproof

    We noted above that $\nil R=\sqrt{\nil R}=\pi^{-1}\nil(R/\nil(R))$.
    But we have also that $\nil R=\pi^{-1}0$, so $\pi^{-1}0=\pi^{-1}\nil(R/\nil(R))$ and since $\pi$ is surjective this means $\nil(R/\nil(R))=0$
    so that $R/\nil(R)$ is reduced.

    Clearly if $x\in\nil(R)$ then $x\in\p$ for every prime $\p$ since $x^n=0\in\p$.
    Now if $x\notin\nil(R)$ then we claim there exists a prime ideal $\p$ such that $x\notin\p$.
    We prove this using Zorn.
    Let $T=\set{1,x,x^2,\dots}$ and let $S$ be the poset of all ideals $I$ such that $I\cap T=\emptyset$; note that all ideals in $S$ are proper.
    As explained before the union of a chain of ideals is an ideal, and clearly if all of the ideals in the union don't intersect $T$, neither
    does the union.
    Thus by Zorn, $S$ has a maximal element $\p$.

    We claim that $\p$ is prime.
    Indeed, suppose that $ab\in\p$ but $a,b\notin\p$.
    Then by maximality we have that $x^n\in(a,\p)$ and $x^m\in(b,\p)$ for some $n,m\geq0$.
    But then $x^{n+m}\in(a,\p)(b,\p)\subseteq\p$, a contradiction.

    Thus there is a prime ideal $\p$ such that $\p\cap T=\emptyset$, in particular $x\notin\p$.
    So if $x\notin\nil(R)$ then $x\notin\bigcap_\p\p$, as desired.
    \qed

\eproof

\bcoro

    Let $I$ be an ideal of $R$, then
    $$ \sqrt I = \bigcap_{I\leq\p}\p , $$
    where the intersection is over all prime ideals containing $I$.

\ecoro

\bproof

    We showed above that $\sqrt I=\pi^{-1}\nil(R/I)$.
    By above $\nil(R/I)$ is the intersection over all prime ideals of $R/I$, which are all prime ideals of $R$ containing $I$:
    $$ \sqrt I = \pi^{-1}\bigcap_{I\leq\p}\bar\p = \bigcap_{I\leq\p}\pi^{-1}\bar\p = \bigcap_{I\leq\p}\p . \qed $$

\eproof

\bdefn

    Let $R$ be a ring, then define its {\emphcolor Jacobson radical} to be
    $$ J(R) = \bigcap_\m\m , $$
    the intersection of all maximal ideals.

\edefn

\blemm

    Let $\m\leq R$ be maximal.
    Then $x\notin\m$ iff there is a $y\in R$ such that $1-xy\in\m$.

\elemm

\bproof

    $x\notin\m$ iff $\bar x\in R/\m$ is nonzero, or equivalently is a unit.
    This is iff there is a $y\in R$ such that $\bar x\bar y=\bar 1$.
    This just means $1-xy\in\m$, as desired.
    \qed

\eproof

\bcoro

    For a ring $R$, its Jacobson radical $J(R)$ is the set of all elements $x\in R$ such that $1-yx\in R^\times$ for all $y\in R$.

\ecoro

\bproof

    Now if $x\notin J(R)$ then $x\notin\m$ for some $\m$, and then $1-yx\in\m$ for some $y$ and thus is not a unit.
    And if $1-yx$ is not a unit, then $1-yx\in\m$ for some maximal $\m$ and thus $x\notin\m$ so that $x\notin J(R)$.
    \qed

\eproof

\bdefn

    A ring $R$ is {\emphcolor local} if it has a unique maximal ideal $\m$.
    The field $k=R/\m$ is called its {\emphcolor residue field}.
    By abuse of notation we write $(R,\m)$ for a local field.

\edefn

\bprop

    \benum
        \item If $(R,\m)$ is a local ring then $R-\m\subseteq R^\times$.
        \item Conversely if $I<R$ is a proper ideal with $R-I\subseteq R^\times$ then $(R,I)$ is a local ring.
        \item If $\m\leq R$ is maximal such that $1+\m\subseteq R^\times$ then $(R,\m)$ is local.
    \eenum

\eprop

\bproof

    If $x\notin\m$ then $(x)$ is not contained in $\m$, and since $\m$ is the unique maximal ideal this means that $(x)=R$ so that $x$ is a unit.
    Also we noted that $R^\times$ is the intersection of all complements of maximal ideals, so this follows from that as well.

    Conversely, we have that
    $$ R-I \subseteq R^\times \bigcap_\m\m^c \implies I \supseteq \bigcup_\m\m , $$
    which means that $I$ contains all maximal ideals, which can only happen if it is the unique maximal ideal.

    We want to show that $R-\m\subseteq R^\times$.
    Indeed, if $x\notin\m$ then there is a $y$ such that $1-yx\in\m$ so that $yx\in1+\m$ and is thus a unit.
    So $x$ is a unit, as desired.
    \qed

\eproof

\bexam

    For $p\in\bZ$ prime, define
    $$ \bZ_{(p)} = \set{\frac ab}[a,b\in\bZ,b\notin(p)] \subseteq \bQ ,\qquad I = \set{\frac ab\in\bZ_{(p)}}[a\in(p)] . $$
    Then $(\bZ_{(p)},I)$ is local.
    Indeed, if $\frac ab\notin I$ then it is invertible (its inverse is $\frac ba$).
    We can generalize this, as we will in the next section.

\eexam

\subsection{Localization of rings}

Consider the following problem: for a ring $A$ and a subset $S\subseteq A$, we want a ring $S^{-1}A$ and a morphism $\iota\colon A\to S^{-1}A$
such that any morphism $f\colon A\to B$ such that $f(S)\subseteq B^\times$, $f$ factors uniquely through $\iota$.
That is, there is a unique $\hat f\colon S^{-1}A\to B$ such that $f=\hat f\circ\iota$:

\centerline{\drawdiagram{
    $A$&$S^{-1}A$\cr
    &$B$
}{
    \diagarrow{from={1,1}, to={1,2}, text=$\iota$, y distance=-.25cm}
    \diagarrow{from={1,1}, to={2,2}, text=$f$, y distance=-.25cm}
    \diagarrow{from={1,2}, to={2,2}, text=$\exists!$, x distance=.25cm, dashed}
}}

This is a universal property, and as such any solution is unique up to unique isomorphism.
In this section we will investigate the construction of such a solution.

\bdefn

    A subset $S\subseteq A$ is {\emphcolor multiplicatively closed} if $1\in S$ and for $a,b\in S$ we have $ab\in S$.

\edefn

It is sufficient to consider only the case that $S$ is multiplicatively closed.
Otherwise, let $\hat S$ be $S$'s multiplicative closure: the set of all finite products of elements from $S$.
Then $\hat S^{-1}R$ is a solution to the problem for $S$.
Indeed, if $f(S)\subseteq B^\times$, then $f(\hat S)\subseteq B^\times$ as well, as the product of units.
Then there is a unique factoring through $\hat S$'s $\iota$.

\bdefn

    For $S\subseteq A$ multiplicatively closed, define an equivalence on $A\times S$ by $(a,s)\sim(b,t)$ iff there is a
    $u\in S$ such that $u(at-bs)=0$, in which case we write $(a,s)\sim_u(b,t)$.

\edefn

\blemm

    This is indeed an equivalence relation.

\elemm

\bproof

    The subrelations $\sim_u$ are all clearly reflexive and symmetric, thus their union $\sim$ is as well.
    Now if $(a,s)\sim_\alpha(b,t)\sim_\beta(c,v)$ then $\alpha at=\alpha sb$ and $\beta bv=\beta ct$.
    So then $\alpha t\beta av=\alpha sb\beta v=\alpha\beta tcs$ so that $(a,s)\sim_{\alpha t\beta}(c,v)$.
    \qed

\eproof

\bexam

    A more naive equivalence would be $\sim_1$, i.e.\ $(a,s)\sim(b,t)$ iff $at-bs=0$.
    This corresponds to $\sim$ when $A$ is an integral domain, but is not generally an equivalence.

\eexam

\bdefn

    For $S\subseteq A$ multiplicatively closed, define $S^{-1}A$ to be the quotient of $A\times S$ by $\sim$.
    Write the equivalence class of $(a,s)$ by $a/s$.
    We give this a ring structure by
    $$ 0 = 0/1,\qquad 1 = 1/1,\qquad a/s + b/t = (at+bs)/st,\qquad a/s \cdot b/t = (ab)/(st) . $$
    Define $\iota\colon A\to S^{-1}A$ by $\iota(a)=a/1$.

\edefn

\bprop

    $S^{-1}A$ is indeed a ring, and $\iota$ is a ring morphism.

\eprop

\bproof

    We need to show that addition and multiplication are well-defined.
    Indeed, if $(a,s)\sim_u(a',s')$ and $(b,t)\sim_v(b',t')$ then $(at+bs,st)\sim_{uv}(a't'+b's',s't')$ which can be verified.
    And similarly $(ab,st)\sim_{uv}(a'b',s't')$.
    From well-definedness it is easy to see that addition and multiplication define a ring structure and $\iota$ is a ring morphism.
    \qed

\eproof

A few properties are immediate:
\benum
    \item For $s\in S$, $\iota(s)=1/s$ has an inverse $s/1$, so that $\iota(S)\subseteq(S^{-1}R)^\times$;
    \item We have that $a/s=0/1$ iff there is a $t\in S$ such that $at=0$.
    In particular, $\iota(a)=0$ iff there is a $t\in S$ such that $ta=0$.
    \item Every element of $S^{-1}A$ is of the form
    $$ a/s = a/1\cdot1/s = \iota(a)\cdot\iota(s)^{-1} . $$
\eenum

\bthrm

    This construction solves the above universal property.

\ethrm

\bproof

    Let $f\colon A\to B$ be such that $f(S)\subseteq B^\times$: we need to show that there is a unique map $g\colon S^{-1}A\to B$
    such that $f=g\circ\iota$.
    Indeed, we must define
    $$ g(a/s) = g(\iota(a)\iota(s)^{-1}) = (g\circ\iota(a))\cdot(g\circ\iota(s))^{-1} = f(a)f(s)^{-1} . $$
    This shows that $g$ is unique, if it exists.
    We need to show that this map is well-defined and a ring morphism.

    If $a/s=b/t$ then we must have that $f(a)f(s)^{-1}=f(b)f(t)^{-1}$ so that $g$ is well-defined.
    Indeed, if $u(at-bv)=0$ for $u\in S$, we have
    $$ f(u)\cdot(f(at)-f(bv)) = 0 $$
    and since $f(S)\subseteq B^\times$ we have that $f(u)$ is a unit so that $f(at)-f(bv)$, meaning this is well-defined.
    And this is a ring morphism, as can be checked easily.
    \qed

\eproof

\bcoro

    Let $f\colon A\to B$ be a ring morphism such that
    \benum
        \item $f(S)\subseteq B^\times$,
        \item if $f(a)=0$ then $ta=0$ for some $t\in S$,
        \item every element of $B$ is of the form $f(a)f(s)^{-1}$ for $a\in A$ and $s\in S$.
    \eenum
    Then there is a unique isomorphism $g\colon S^{-1}A\to B$ such that $g\circ\iota=f$.

\ecoro

\bproof

    $f$ solves the universal problem above for $S$ as well.
    If $h\colon A\to C$ is a map where $h(S)\subseteq C^\times$, then we need to define $\hat h\colon B\to C$ by
    $$ \hat h(f(a)f(s)^{-1}) = h(a)h(s)^{-1} , $$
    which is well-defined and a ring morphism for similar reasons as above.
    Since universal properties are unique up to unique isomorphism, the result follows.

    Alternatively we can prove this directly: let $g\colon S^{-1}A\to B$ be the unique map where $g\circ\iota=f$ by the universal property.
    By definition $g(a/s)=f(a)f(s^{-1})$.
    By the third point $g$ is surjective, and it is injective since if $g(a/s)=0$ we have $f(a)=0$ so that $ta=0$, which means that $a/s=0$.
    \qed

\eproof

\bexam

    If $A$ is an integral domain, then $S=A-0$ is multiplicatively closed, and $S^{-1}A$ is the field of fractions of $A$.

\eexam

\bprop

    For $A\neq0$, then $S^{-1}A=0$ iff $0\in S$.

\eprop

\bproof

    We have that $S^{-1}A=0$ iff $1/1=0/1$ iff there is a $s\in S$ such that $s\cdot1=0$.
    \qed

\eproof

\bdefn

    Let $\p\subseteq A$ be prime.
    Then $S=A-\p$ is multiplicatively closed, and we denote its localization by $A_\p=S^{-1}A$.

    For $a\in A$, $S=\set{1,a,a^2,\dots}$ is multiplicatively closed, and we denote its localization by $A_a=S^{-1}A$.

\edefn

\bnote

    Note that we have two distinct localizations here, which have similar names in some cases.
    If $A$ is a PID and $p\in A$ is irreducible, then $A_{(p)}$ is localization by $A-(p)$, which $A_p$ is localization by $\set{p^n}_{n<\omega}$.
    Elements of the former are of the form $a/s$ where $p$ does not divide $s$, while elements of the former are of the form $a/s$ where $s=p^n$.
    Quite different!

\enote

\bdefn

    For a ring morphism $f\colon A\to B$ and an ideal $I\subseteq A$, define the ideal $IB=(f(I))$ of $B$.
    Explicitly, $IB$ is the set of all finite sums $\sum_if(a_i)b_i$ for $a_i\in I$.

\edefn

In general we have $f^{-1}(f(I))\supseteq I$ and $f(f^{-1}(J))\subseteq J$.
These imply $f^{-1}(IB)\supseteq I$ and $f^{-1}(J)B\subseteq J$.

\bdefn

    For two ideals $I,J$, define $(I:J)=\set{x\in A}[xJ\subseteq I]$.
    This is an ideal containing $I$.
    For $a\in A$, define $(I:a)=(I:(a))=\set{x\in A}[xa\subseteq I]$.

    For $S\subseteq R$ multiplicatively closed, define the ideal $S^{-1}I=IS^{-1}A\subseteq S^{-1}A$.
    This is the ideal generated by $\iota(I)$.

\edefn

So we have two maps: $S^{-1}\bul$ mapping ideals of $A$ to ideals of $S^{-1}A$, and $\iota^{-1}$ mapping ideals of $S^{-1}A$ to $A$:
$$ \set{\hbox{Ideals of $A$}} \xtof{\quad S^{-1}\quad}[\quad\iota^{-1}\quad] \set{\hbox{Ideals of $S^{-1}A$}} . $$
We wish to relate the two.

\bprop

    \benum
        \item For an ideal $I\subseteq A$, $S^{-1}I=\set{a/s}[a\in I,s\in S]$.
        \item For an ideal $J\subseteq S^{-1}A$, we have $S^{-1}\iota^{-1}J=J$.
        \item For an ideal $I\subseteq A$, we have $\iota^{-1}S^{-1}I=\bigcup_{s\in S}(I:s)$.
        \item We have $S^{-1}I=S^{-1}A$ iff $I\cap S\neq\emptyset$.
        \item For $\p\subseteq A$ prime, if $I\cap S\neq\emptyset$ then $S^{-1}\p\subseteq S^{-1}A$ is prime (otherwise it is all of $S^{-1}A$).
        \item The correspondence $\q\mapsto\iota^{-1}\q$ is a bijection between prime ideals of $S^{-1}A$ and prime ideals of $A$ not intersecting $S$.
        Its inverse is $S^{-1}$.
        That is, we have a one-to-one correspondence:
        $$ \set{\hbox{Prime ideals of $A$ not intersecting $S$}} \xtof{\quad S^{-1}\quad}[\quad\iota^{-1}\quad] \set{\hbox{Prime ideals of $S^{-1}A$}} . $$
    \eenum

\eprop

\bproof

    \benum
        \item $\set{a/s}[a\in I,s\in S]$ is an ideal and contains $\iota(I)$.
        It is clearly closed under multiplication by $S^{-1}A$ and for $a/s+b/t=(at+bs)/(st)$, notice that $at+bs\in I$.
        By definition $S^{-1}I$ is the smallest ideal containing $\iota(I)$, and it consists of finite sums $\sum_i\iota(a_i)x_i/s_i$.
        This clearly contains $a/s$, and thus they must be equal.

        \item As explained before, in general we have $S^{-1}\iota^{-1}J=\iota^{-1}(J)S^{-1}A\subseteq J$, we must prove the opposite inclusion.
        As above, we have $S^{-1}\iota^{-1}J=\set{a/s}[\iota(a)\in J,s\in S]$.
        For $x/s\in J$, we have that $\iota(x)=x/1\in J$ and so $x/s\in S^{-1}\iota^{-1}J$ as desired.

        \item If $a\in\iota^{-1}S^{-1}I$ then $\iota(a)\in S^{-1}I$ so by the first point $a/1=b/s$ for $b\in I$.
        Then $asu=bu$ for some $u\in S$ and since $bu\in J$ this means that $a\in(I:su)$.
        Conversely if $a\in(I:s)$ for $s\in S$ then $as\in I$ and so $\iota(a)=a/1=(as)/s\in S^{-1}I$.

        \item If $S^{-1}I=S^{-1}A$ then $1/s=a/t$ for $a\in I$, so $as=t$ meaning $as\in I\cap S$.
        Conversely if $a\in I\cap S$ then $1/1=a/a\in S^{-1}I$.

        \item If $S\cap\p=\emptyset$ then we need to show that $S^{-1}\p$ is prime.
        If $a_1/s_1,a_2/s_2\notin S^{-1}\p$ we want to show $(a_1a_2)/(s_1s_2)\notin S^{-1}\p$.
        By assumption $a_1,a_2\notin\p$.

        Now we claim that if $a\notin\p$, $a/s\notin S^{-1}\p$, which would complete the proof.
        Assume the contrary: $a/s=b/t$ for $b\in\p$ then there is a $u\in S$ such that $atu=bsu$.
        Then $atu\in\p$, but $a\notin\p$ and $t,u\in S$ and thus also are not in $\p$, a contradiction to $\p$'s primality.

        \item Preimages of prime ideals are prime, so $\iota^{-1}$ maps prime ideals to prime ideals.
        By the second point $S^{-1}\iota^{-1}\q=\q$ for all prime ideals $\q\subseteq S^{-1}A$.
        It remains to show that $\iota^{-1}S^{-1}\p=\p$ for all prime $\p\subseteq A$ which don't intersect $S$.
        We know that $\p\subseteq\iota^{-1}S^{-1}\p$ as is true in general, so we want to show that $\iota^{-1}S^{-1}\p=\bigcup_{s\in S}(\p:s)$ is contained in $\p$.
        Indeed, if $a\in(\p:s)$ then $as\in\p$ and since $s\notin\p$ we have that $a\in\p$ as desired.
        \qed
    \eenum

\eproof

\bcoro

    Let $\p\subseteq A$ be prime, and consider its localization $A_\p$.
    Then $A_\p$ is a local ring with maximal ideal $\p A_\p$.
    Moreover, the correspondence $\q\to\q A_\p$ is a bijection between prime ideals of $A_\p$ and prime ideals $\q\subseteq\p$.

\ecoro

\bproof

    Let $S=A-\p$, so that $A_\p=S^{-1}A$, and $\p A_\p=S^{-1}\p$.
    The last point is clear: $\q\cap S=\emptyset$ iff $\q\subseteq\p$.
    $A_\p$ being local is also clear: $S^{-1}\bul$ is order-preserving, so $S^{-1}\q\subseteq S^{-1}\p$ for all prime ideals $\q$.
    Thus $S^{-1}\p$ is maximal.
    \qed

\eproof

\bcoro

    If $a\in A$ is not nilpotent, then there is a prime ideal $\p\subseteq A$ such that $a\notin\p$.

\ecoro

\bproof

    As before, consider the multiplicative set $S=\set{a^n}_{n<\omega}$, and $A_a=S^{-1}A$.
    Since $a$ is not nilpotent, $0\notin S$, so that $A_a\neq0$.
    So take a prime ideal $\p\subseteq A_a$, and then $\iota^{-1}\p$ is a prime ideal of $A$ not intersecting $S$, thus not containing $a$.
    \qed

\eproof

\bnote

    This completes the proof before that
    $$ \nil(A) = \bigcap_\p\p $$
    but without appealing to Zorn, and thus not relying on the Axiom of Choice.

\enote

